% Options for packages loaded elsewhere
\PassOptionsToPackage{unicode}{hyperref}
\PassOptionsToPackage{hyphens}{url}
\PassOptionsToPackage{dvipsnames,svgnames,x11names}{xcolor}
%
\documentclass[
  11pt,
]{article}
\usepackage{amsmath,amssymb}
\usepackage{iftex}
\ifPDFTeX
  \usepackage[T1]{fontenc}
  \usepackage[utf8]{inputenc}
  \usepackage{textcomp} % provide euro and other symbols
\else % if luatex or xetex
  \usepackage{unicode-math} % this also loads fontspec
  \defaultfontfeatures{Scale=MatchLowercase}
  \defaultfontfeatures[\rmfamily]{Ligatures=TeX,Scale=1}
\fi
\usepackage{lmodern}
\ifPDFTeX\else
  % xetex/luatex font selection
  \setmainfont[]{DejaVu Serif}
  \setmonofont[Scale=0.82]{DejaVu Sans Mono}
\fi
% Use upquote if available, for straight quotes in verbatim environments
\IfFileExists{upquote.sty}{\usepackage{upquote}}{}
\IfFileExists{microtype.sty}{% use microtype if available
  \usepackage[]{microtype}
  \UseMicrotypeSet[protrusion]{basicmath} % disable protrusion for tt fonts
}{}
\makeatletter
\@ifundefined{KOMAClassName}{% if non-KOMA class
  \IfFileExists{parskip.sty}{%
    \usepackage{parskip}
  }{% else
    \setlength{\parindent}{0pt}
    \setlength{\parskip}{6pt plus 2pt minus 1pt}}
}{% if KOMA class
  \KOMAoptions{parskip=half}}
\makeatother
\usepackage{xcolor}
\usepackage[top=2.5cm,bottom=2.5cm,left=2.8cm,right=2.8cm]{geometry}
\usepackage{listings}
\newcommand{\passthrough}[1]{#1}
\lstset{defaultdialect=[5.3]Lua}
\lstset{defaultdialect=[x86masm]Assembler}
\usepackage{longtable,booktabs,array}
\usepackage{calc} % for calculating minipage widths
% Correct order of tables after \paragraph or \subparagraph
\usepackage{etoolbox}
\makeatletter
\patchcmd\longtable{\par}{\if@noskipsec\mbox{}\fi\par}{}{}
\makeatother
% Allow footnotes in longtable head/foot
\IfFileExists{footnotehyper.sty}{\usepackage{footnotehyper}}{\usepackage{footnote}}
\makesavenoteenv{longtable}
\setlength{\emergencystretch}{3em} % prevent overfull lines
\providecommand{\tightlist}{%
  \setlength{\itemsep}{0pt}\setlength{\parskip}{0pt}}
\setcounter{secnumdepth}{-\maxdimen} % remove section numbering
\ifLuaTeX
  \usepackage{selnolig}  % disable illegal ligatures
\fi
\IfFileExists{bookmark.sty}{\usepackage{bookmark}}{\usepackage{hyperref}}
\IfFileExists{xurl.sty}{\usepackage{xurl}}{} % add URL line breaks if available
\urlstyle{same}
\hypersetup{
  colorlinks=true,
  linkcolor={blue},
  filecolor={Maroon},
  citecolor={Blue},
  urlcolor={Blue},
  pdfcreator={LaTeX via pandoc}}

\title{SpeakPrep: Product Requirements Document}
\author{Abhiyan Sainju}
\date{April 2026}

\begin{document}
\maketitle

\hypertarget{speakprep-document-1-product-requirements-document}{%
\section{SpeakPrep --- Document 1: Product Requirements
Document}\label{speakprep-document-1-product-requirements-document}}

\hypertarget{version-1.0-author-abhiyan-sainju-april-2026}{%
\subsubsection{Version 1.0 \textbar{} Author: Abhiyan Sainju \textbar{}
April 2026}\label{version-1.0-author-abhiyan-sainju-april-2026}}

\hypertarget{status-living-document-update-as-product-evolves}{%
\subsubsection{Status: LIVING DOCUMENT --- update as product
evolves}\label{status-living-document-update-as-product-evolves}}

\begin{center}\rule{0.5\linewidth}{0.5pt}\end{center}

\begin{quote}
\textbf{How to use this doc:} This is your single source of truth on
WHAT you're building and WHY. Every engineering decision in Docs 2 and 3
traces back here. When you're mid-build and tempted to add a feature,
come back to this doc first. If it doesn't serve a user story here, it
doesn't get built yet.
\end{quote}

\begin{center}\rule{0.5\linewidth}{0.5pt}\end{center}

\hypertarget{table-of-contents}{%
\subsection{TABLE OF CONTENTS}\label{table-of-contents}}

\begin{enumerate}
\def\labelenumi{\arabic{enumi}.}
\tightlist
\item
  \protect\hyperlink{1-executive-summary}{Executive Summary}
\item
  \protect\hyperlink{2-problem-statement}{Problem Statement}
\item
  \protect\hyperlink{3-market-context}{Market Context}
\item
  \protect\hyperlink{4-target-users}{Target Users}
\item
  \protect\hyperlink{5-product-vision}{Product Vision}
\item
  \protect\hyperlink{6-feature-set--phase-1}{Feature Set --- Phase 1
  (Interview Coach)}
\item
  \protect\hyperlink{7-feature-set--phase-2-deferred}{Feature Set ---
  Phase 2 (General Platform)}
\item
  \protect\hyperlink{8-user-stories}{User Stories}
\item
  \protect\hyperlink{9-non-goals}{Non-Goals}
\item
  \protect\hyperlink{10-success-metrics--kpis}{Success Metrics \& KPIs}
\item
  \protect\hyperlink{11-assumptions-vs-decisions}{Assumptions vs
  Decisions}
\item
  \protect\hyperlink{12-risks}{Risks}
\item
  \protect\hyperlink{13-monetization-strategy}{Monetization Strategy}
\item
  \protect\hyperlink{14-competitive-differentiation}{Competitive
  Differentiation}
\end{enumerate}

\begin{center}\rule{0.5\linewidth}{0.5pt}\end{center}

\hypertarget{executive-summary}{%
\subsection{1. Executive Summary}\label{executive-summary}}

\textbf{SpeakPrep} is a real-time, voice-first AI mock interview coach.
Users speak naturally; SpeakPrep responds conversationally --- probing,
pushing back, and scoring --- in under 2 seconds. It covers the full
interview lifecycle: resume intake, behavioral practice, technical
walkthroughs, and system design sessions, with AI feedback that improves
with every session.

SpeakPrep is \textbf{Phase 1} of a larger voice AI platform. The
interview coach is the wedge --- a high-urgency, high-retention use case
that proves the core voice pipeline. The platform expands to general
voice AI in Phase 2.

\textbf{Why interview coaching as the wedge:} - High personal urgency
(users are actively job-hunting --- pain is immediate) - Repeated usage
need (multiple interview rounds, multiple companies) - Measurable
outcomes (did they get the job? did their scores improve?) - Defensible
against ChatGPT/Claude (those are chat products; we're a voice
\emph{coach}) - Personal relevance: this is your own lived experience as
a recent MS grad job hunting

\begin{center}\rule{0.5\linewidth}{0.5pt}\end{center}

\hypertarget{problem-statement}{%
\subsection{2. Problem Statement}\label{problem-statement}}

\hypertarget{the-primary-problem}{%
\subsubsection{The Primary Problem}\label{the-primary-problem}}

Job seekers need realistic, conversational interview practice. But:

\textbf{Human mocks are inaccessible.} Professional coaching costs
\$150--\$500/session. Peer mocks on Pramp/Exponent yield inconsistent
quality --- 30\% of users report bad experiences. Family/friends can't
replicate interviewer pressure.

\textbf{Current AI tools don't feel like real interviews.} Google
Interview Warmup asks questions and waits. Final Round AI reads back
bullet points. ChatGPT doesn't push back, probe for details, or simulate
an interviewer who says ``walk me through your algorithm choice.'' None
of these products create the \emph{social pressure} that makes real
interviews hard.

\textbf{The prep journey is fragmented.} A typical job seeker uses 4--5
separate tools: LeetCode for coding (\$35/month), Big Interview for
behavioral (\$39--\$299/month), Glassdoor for company research (free),
Interviewing.io for calibration (\$225/session), ChatGPT for stories (ad
hoc). No product covers the whole journey.

\textbf{Progress is invisible.} Users practice but don't know if they're
improving. There's no score trend, no ``you ramble on this topic,'' no
``your STAR answers are missing Results 60\% of the time.''

\hypertarget{the-secondary-problem-platform-vision}{%
\subsubsection{The Secondary Problem (Platform
Vision)}\label{the-secondary-problem-platform-vision}}

Beyond interviews, voice interaction is fundamentally better than text
for many AI use cases --- language practice, presentation coaching,
thinking out loud, real-time tutoring. Building the voice pipeline for
SpeakPrep creates the infrastructure for these adjacent products.

\begin{center}\rule{0.5\linewidth}{0.5pt}\end{center}

\hypertarget{market-context}{%
\subsection{3. Market Context}\label{market-context}}

\hypertarget{market-size}{%
\subsubsection{Market Size}\label{market-size}}

\begin{itemize}
\tightlist
\item
  AI recruitment market: \textbf{\$596--706M in 2025}, growing at
  6.8--7.5\% CAGR
\item
  Candidate-facing AI interview coaching segment: \textbf{\$1.7B in
  2025}, growing at \textbf{30.6\% CAGR}
\item
  58\% of job applicants now use AI tools during their search
\item
  93\% of job seekers report interview anxiety as a significant barrier
\end{itemize}

\hypertarget{competitive-landscape}{%
\subsubsection{Competitive Landscape}\label{competitive-landscape}}

\begin{longtable}[]{@{}
  >{\raggedright\arraybackslash}p{(\columnwidth - 6\tabcolsep) * \real{0.2500}}
  >{\raggedright\arraybackslash}p{(\columnwidth - 6\tabcolsep) * \real{0.2500}}
  >{\raggedright\arraybackslash}p{(\columnwidth - 6\tabcolsep) * \real{0.2500}}
  >{\raggedright\arraybackslash}p{(\columnwidth - 6\tabcolsep) * \real{0.2500}}@{}}
\toprule\noalign{}
\begin{minipage}[b]{\linewidth}\raggedright
Competitor
\end{minipage} & \begin{minipage}[b]{\linewidth}\raggedright
Price
\end{minipage} & \begin{minipage}[b]{\linewidth}\raggedright
What They Do
\end{minipage} & \begin{minipage}[b]{\linewidth}\raggedright
Fatal Weakness
\end{minipage} \\
\midrule\noalign{}
\endhead
\bottomrule\noalign{}
\endlastfoot
\textbf{Final Round AI} & \$96--\$148/mo & Live CoPilot during
interviews, AI mock & Buggy, expensive, generic responses. Users report
it ``feels like a spell-checker, not a coach.'' \\
\textbf{Google Interview Warmup} & Free & 5 scripted questions, basic
keyword feedback & No conversation, no follow-ups, no personalization.
Completely static. \\
\textbf{Yoodli} & \$10--\$20/mo & Speech delivery coaching (pacing,
fillers) & Coaches delivery only, not content. Won't tell you your STAR
answer was weak. \\
\textbf{Hume AI} & API only & Voice AI infrastructure with emotional
detection & Not a consumer product. Requires engineering to use. \\
\textbf{Interviewing.io} & \$225+/session & Live mock with FAANG
engineers & \$1K+ for 5 sessions. Inaccessible to most job seekers. \\
\textbf{Pramp/Exponent} & Free--\$79/mo & Peer-to-peer mocks & 30\% bad
peer experiences. No consistency. \\
\textbf{Big Interview} & \$39--\$299 & Video lessons + scripted question
practice & Static, not conversational. AI feedback is generic. \\
\end{longtable}

\hypertarget{the-gap-speakprep-fills}{%
\subsubsection{The Gap SpeakPrep Fills}\label{the-gap-speakprep-fills}}

No competitor combines all four of: 1. \textbf{Real-time voice
conversation} (not text, not scripted Q\&A) 2. \textbf{Intelligent
probing} (follow-up questions based on your actual answer) 3.
\textbf{End-to-end coverage} (behavioral + technical + system design in
one product) 4. \textbf{Progress tracking} (scores trend over time,
specific improvement areas)

\begin{center}\rule{0.5\linewidth}{0.5pt}\end{center}

\hypertarget{target-users}{%
\subsection{4. Target Users}\label{target-users}}

\hypertarget{primary-user-the-active-job-seeker-jane}{%
\subsubsection{Primary User: ``The Active Job Seeker''
(Jane)}\label{primary-user-the-active-job-seeker-jane}}

\textbf{Who she is:} - Age 22--32, recent graduate or early-career
professional (0--5 YOE) - Actively applying for software engineering or
data roles - Has 3--8 interviews scheduled or expected in the next 4
weeks - Has tried ChatGPT for practice but finds it too passive - Can't
afford \$150/session human coaching

\textbf{Her pain:} \textgreater{} ``I know what I want to say but when I
actually have to say it out loud in an interview my mind goes blank. I
need to practice \emph{speaking}, not just writing.''

\textbf{What she needs:} - A practice partner that pushes back like a
real interviewer - Specific feedback (``your Result was vague --- what
was the actual business impact?'') - Confidence from repetition before
the real thing

\textbf{Session pattern:} 30--45 min/day, 3--5x per week, for 2--6 weeks
before interviews.

\begin{center}\rule{0.5\linewidth}{0.5pt}\end{center}

\hypertarget{secondary-user-the-career-switcher-marcus}{%
\subsubsection{Secondary User: ``The Career Switcher''
(Marcus)}\label{secondary-user-the-career-switcher-marcus}}

\textbf{Who he is:} - Age 28--40, transitioning from non-technical
background (finance, law, sales) - Pursuing product management, data
analytics, or SWE roles - Less interview experience, more anxiety -
Values structured guidance over open-ended practice

\textbf{His pain:} \textgreater{} ``I don't even know what interviewers
are looking for. I need someone to explain what a good answer sounds
like, not just ask me questions.''

\textbf{What he needs:} - More educational scaffolding (explain what
STAR means, show examples) - Gentler difficulty progression -
Post-session written feedback to study

\begin{center}\rule{0.5\linewidth}{0.5pt}\end{center}

\hypertarget{tertiary-user-the-non-native-speaker-wei}{%
\subsubsection{Tertiary User: ``The Non-Native Speaker''
(Wei)}\label{tertiary-user-the-non-native-speaker-wei}}

\textbf{Who he is:} - International student or professional, strong
technical skills - English is second language --- confident in writing,
less confident speaking - Particularly anxious about phone screens and
verbal communication

\textbf{His pain:} \textgreater{} ``I answer technical questions fine
but when they ask `tell me about yourself' I freeze. I need to practice
speaking English in an interview context specifically.''

\textbf{What he needs:} - Speaking pace and clarity feedback -
Pronunciation-aware scoring (not penalize accents, but flag pace/filler
issues) - Low-stakes environment to build spoken English confidence

\begin{center}\rule{0.5\linewidth}{0.5pt}\end{center}

\hypertarget{product-vision}{%
\subsection{5. Product Vision}\label{product-vision}}

\hypertarget{phase-1-product-speakprep-interview-coach-what-were-building-now}{%
\subsubsection{Phase 1 Product (SpeakPrep Interview Coach) --- What
We're Building
Now}\label{phase-1-product-speakprep-interview-coach-what-were-building-now}}

A web-based, voice-first AI mock interview platform where:

\begin{enumerate}
\def\labelenumi{\arabic{enumi}.}
\tightlist
\item
  \textbf{User uploads resume + selects target role/company}
\item
  \textbf{Selects interview type:} Behavioral, Technical, System Design,
  Full Loop
\item
  \textbf{Speaks naturally} --- SpeakPrep transcribes, understands,
  responds in real-time
\item
  \textbf{AI probes intelligently} --- ``Can you be more specific about
  what you personally contributed?'' --- not scripted follow-ups
\item
  \textbf{Scoring runs live and post-session} --- content, structure,
  confidence markers, filler words
\item
  \textbf{Post-session report} --- per-question breakdown, trends,
  specific improvement actions
\item
  \textbf{Progress tracked across sessions} --- are you actually getting
  better?
\end{enumerate}

\hypertarget{phase-2-product-general-voice-ai-platform-deferred}{%
\subsubsection{Phase 2 Product (General Voice AI Platform) ---
Deferred}\label{phase-2-product-general-voice-ai-platform-deferred}}

Once the voice pipeline is proven and user base established: -
\textbf{Speaking coach} (presentations, speeches, public speaking
practice) - \textbf{Language learning} (conversational practice in
target language) - \textbf{Thinking out loud} (rubber duck debugging
with a voice AI) - \textbf{General assistant} (the full-breadth voice AI
use case)

This platform expansion uses identical infrastructure. Only the
prompting, session logic, and UI differ. This is the architectural
decision that makes the general platform achievable: we build it right
once.

\begin{center}\rule{0.5\linewidth}{0.5pt}\end{center}

\hypertarget{feature-set-phase-1}{%
\subsection{6. Feature Set --- Phase 1}\label{feature-set-phase-1}}

\hypertarget{f-001-onboarding-resume-intake}{%
\subsubsection{F-001: Onboarding \& Resume
Intake}\label{f-001-onboarding-resume-intake}}

\textbf{What it does:} - User creates account (email/password or Google
OAuth) - Uploads resume (PDF) --- extracted and stored as structured
data - Sets target: job title, target companies (optional), experience
level - Short calibration (3 questions to baseline current skill level)

\textbf{Why it matters:} Personalized interviews require knowing who the
user is. A generic ``Tell me about a time you showed leadership''
becomes ``You managed a team at ECS Tech in Kathmandu --- walk me
through a leadership challenge from that role'' with resume context.
That specificity is what no current tool has.

\textbf{Technical requirements:} - PDF upload (max 5 MB) → PyMuPDF text
extraction → LLM structured parsing - Pydantic schema: name,
experience{[}{]}, education{[}{]}, skills{[}{]}, projects{[}{]} - Stored
in PostgreSQL JSONB column for flexible querying - Extraction time
target: \textless{} 10 seconds

\begin{center}\rule{0.5\linewidth}{0.5pt}\end{center}

\hypertarget{f-002-session-setup}{%
\subsubsection{F-002: Session Setup}\label{f-002-session-setup}}

\textbf{What it does:} - Choose interview type: Behavioral / Technical /
System Design / Mixed - Choose difficulty: Beginner / Intermediate /
Advanced (or let AI calibrate) - Choose interviewer persona: Friendly /
Neutral / Challenging / Adversarial - Choose duration: 15 min / 30 min /
45 min - Optional: paste job description for company-specific tailoring

\textbf{Why it matters:} Difficulty and persona selection creates the
``social pressure'' simulation that no competitor offers. An adversarial
interviewer who interrupts and pushes back is a fundamentally different
experience than a friendly one --- both are necessary preparation.

\begin{center}\rule{0.5\linewidth}{0.5pt}\end{center}

\hypertarget{f-003-real-time-voice-session}{%
\subsubsection{F-003: Real-Time Voice
Session}\label{f-003-real-time-voice-session}}

\textbf{What it does:} - Push-to-talk or continuous voice input (user
configurable) - Real-time transcription displayed as user speaks
(confidence signal) - AI response in \textless{} 2 seconds end-to-end
(speech → speech) - AI voice: select from 4--6 natural-sounding voices
(Kokoro TTS) - Barge-in: user can interrupt AI mid-response - Session
timer visible, question counter visible - Live filler word counter
(``you've said `um' 4 times this question'')

\textbf{Technical requirements:} - WebSocket connection maintained for
full session duration - Silero VAD for end-of-utterance detection -
Deepgram streaming ASR (primary), Faster-Whisper (fallback) - Groq Llama
70B (primary LLM), Cerebras/Gemini Flash (fallback) - Kokoro TTS for
response audio - All audio: 16 kHz mono PCM → 32 kbps Opus over
WebSocket

\begin{center}\rule{0.5\linewidth}{0.5pt}\end{center}

\hypertarget{f-004-ai-interviewer-intelligence}{%
\subsubsection{F-004: AI Interviewer
Intelligence}\label{f-004-ai-interviewer-intelligence}}

\textbf{What it does:}

The AI doesn't just ask scripted questions. It: - \textbf{Asks
intelligent follow-ups} based on what you actually said (``You mentioned
you `worked on' the migration --- were you the lead or a contributor?'')
- \textbf{Probes for STAR components} when missing (``What was the
actual outcome? Did it ship?'') - \textbf{Challenges weak answers} when
in Challenging/Adversarial mode (``That sounds like what the team
decided --- what would you have done differently?'') -
\textbf{Transitions naturally} between questions without breaking
conversational flow - \textbf{Tracks question coverage} --- doesn't ask
the same topic twice - \textbf{Adapts difficulty} within session based
on answer quality

\textbf{Why it matters:} This is the core product moat. Any competitor
can ask ``Tell me about a time you showed leadership.'' Only SpeakPrep
asks ``You said you `helped coordinate' --- can you tell me what you
specifically did versus what others did? What was your decision
authority?'' That's what a real interviewer does.

\begin{center}\rule{0.5\linewidth}{0.5pt}\end{center}

\hypertarget{f-005-live-scoring-during-session}{%
\subsubsection{F-005: Live Scoring (during
session)}\label{f-005-live-scoring-during-session}}

\textbf{What it displays:} - Filler word count per question (live
ticker) - Response time per question - Current question score estimate
(updates after each answer) - Speaking pace indicator (WPM, target:
140--160)

\textbf{Why it matters:} Real-time feedback changes behavior
immediately. Seeing ``you've said `um' 7 times'' mid-question triggers
self-correction in a way post-session feedback cannot.

\begin{center}\rule{0.5\linewidth}{0.5pt}\end{center}

\hypertarget{f-006-post-session-report}{%
\subsubsection{F-006: Post-Session
Report}\label{f-006-post-session-report}}

\textbf{What it contains:} - Overall session score (0--100) -
Per-question scores across 5 dimensions: - Content Quality (30\%) ---
relevance, depth, accuracy - Communication (25\%) --- clarity,
structure, conciseness - STAR Compliance (20\%) --- all 4 components,
specificity, results - Confidence Markers (15\%) --- assertive language,
hedging frequency - Filler Words (10\%) --- rate per minute - Written
feedback per question: what you did well, what to improve, example of
better answer - Audio playback of your answer (optional, for
self-review) - Transcript of full session

\textbf{Why it matters:} The report is what users share, save, and
reference. It's the ``receipt'' of their practice. Without a detailed
report, sessions feel disposable.

\begin{center}\rule{0.5\linewidth}{0.5pt}\end{center}

\hypertarget{f-007-progress-dashboard}{%
\subsubsection{F-007: Progress
Dashboard}\label{f-007-progress-dashboard}}

\textbf{What it shows:} - Score trend over time (line chart per
dimension) - Total practice time, total questions answered - Strongest
and weakest dimensions (rolling 30-day) - Filler word frequency trend -
Response time trend - ``Your STAR compliance has improved 23\% over the
last 10 sessions'' - Question bank coverage: which topics have you
practiced, which haven't you touched?

\textbf{Why it matters:} Progress visibility is the primary retention
mechanism. Users who see improvement continue. Users who feel stuck
without evidence quit. The dashboard converts single-session users to
long-term users.

\begin{center}\rule{0.5\linewidth}{0.5pt}\end{center}

\hypertarget{f-008-question-bank}{%
\subsubsection{F-008: Question Bank}\label{f-008-question-bank}}

\textbf{What it contains:} - 500+ curated questions across categories: -
Behavioral: Leadership, Teamwork, Conflict, Failure, Decision-Making,
Motivation - Technical: Algorithms, Data Structures, System Concepts,
Language-Specific - System Design: Common systems (URL shortener, rate
limiter, chat system, etc.) - Company-Specific: Common patterns from
Glassdoor/Blind for FAANG, startups - Tagged by: difficulty, category,
subcategory, company, role level - User-contributed questions (Phase 1.5
--- after initial launch)

\textbf{Adaptive selection:} Questions selected based on: - User's skill
gaps from previous sessions - Target company/role context - ELO-based
difficulty targeting 60--75\% expected success rate - Coverage balance
(don't over-index on one category)

\begin{center}\rule{0.5\linewidth}{0.5pt}\end{center}

\hypertarget{f-009-authentication-user-accounts}{%
\subsubsection{F-009: Authentication \& User
Accounts}\label{f-009-authentication-user-accounts}}

\textbf{What it does:} - Email/password registration with email
verification - Google OAuth (via Supabase Auth) - JWT access tokens
(15-min) + refresh tokens (7-day httpOnly cookies) - Protected routes on
frontend - Account settings: email, password change, notification
preferences - Data export (GDPR compliance) and account deletion

\begin{center}\rule{0.5\linewidth}{0.5pt}\end{center}

\hypertarget{f-010-rate-limiting-usage-tiers}{%
\subsubsection{F-010: Rate Limiting \& Usage
Tiers}\label{f-010-rate-limiting-usage-tiers}}

\textbf{Free Tier:} - 3 sessions/week (15 min each) - Basic post-session
report (no audio playback, no full transcript) - 7-day score history

\textbf{Pro Tier (\$12/month):} - Unlimited sessions - Full post-session
reports with audio playback - 90-day score history with trend charts -
All interviewer personas and difficulty levels - Company-specific
question sets - Resume-personalized questions

\begin{center}\rule{0.5\linewidth}{0.5pt}\end{center}

\hypertarget{feature-set-phase-2-deferred}{%
\subsection{7. Feature Set --- Phase 2
(Deferred)}\label{feature-set-phase-2-deferred}}

These features are explicitly \textbf{NOT built in Phase 1}. Document
them here so the architecture accommodates them without rework.

\begin{longtable}[]{@{}
  >{\raggedright\arraybackslash}p{(\columnwidth - 4\tabcolsep) * \real{0.3333}}
  >{\raggedright\arraybackslash}p{(\columnwidth - 4\tabcolsep) * \real{0.3333}}
  >{\raggedright\arraybackslash}p{(\columnwidth - 4\tabcolsep) * \real{0.3333}}@{}}
\toprule\noalign{}
\begin{minipage}[b]{\linewidth}\raggedright
Feature
\end{minipage} & \begin{minipage}[b]{\linewidth}\raggedright
Why Deferred
\end{minipage} & \begin{minipage}[b]{\linewidth}\raggedright
When to Build
\end{minipage} \\
\midrule\noalign{}
\endhead
\bottomrule\noalign{}
\endlastfoot
Mobile app (iOS/Android) & Audio APIs on React Native are unreliable;
4--6 weeks of bugs & After web is stable with 100+ DAU \\
Voice cloning / custom voices & Kokoro doesn't support it; ElevenLabs
costs money & Phase 2 premium feature \\
Multi-language support & ASR accuracy drops on non-English; LLM prompts
need rewrite & Phase 2 after English is solid \\
Live interview CoPilot & Different product (listen to real interview,
whisper answers) & Phase 2 --- raises ethical questions too \\
Speaking coach mode & Same pipeline, different prompts & Phase 2 with
existing infrastructure \\
General voice assistant mode & Same pipeline, different prompts & Phase
2 after core is proven \\
Team/enterprise accounts & B2B sales motion, different product & Phase
3 \\
Video analysis (body language) & Completely different tech stack (CV
models) & Long-term \\
\end{longtable}

\begin{center}\rule{0.5\linewidth}{0.5pt}\end{center}

\hypertarget{user-stories}{%
\subsection{8. User Stories}\label{user-stories}}

Organized by feature. Format:
\passthrough{\lstinline!US-[ID]: As a [user], I want to [action], so that [outcome].!}

\hypertarget{onboarding}{%
\subsubsection{Onboarding}\label{onboarding}}

\begin{itemize}
\tightlist
\item
  \passthrough{\lstinline!US-001!} As a new user, I want to upload my
  resume so that interview questions are personalized to my experience.
\item
  \passthrough{\lstinline!US-002!} As a new user, I want to set my
  target role and company so that questions are calibrated to what I'm
  actually preparing for.
\item
  \passthrough{\lstinline!US-003!} As a new user, I want a calibration
  session so that the AI doesn't start me too easy or too hard.
\end{itemize}

\hypertarget{session}{%
\subsubsection{Session}\label{session}}

\begin{itemize}
\tightlist
\item
  \passthrough{\lstinline!US-010!} As a user, I want to press a button
  and speak naturally so that I can practice without typing.
\item
  \passthrough{\lstinline!US-011!} As a user, I want to hear the AI
  respond in a natural voice within 2 seconds so that the conversation
  feels real.
\item
  \passthrough{\lstinline!US-012!} As a user, I want the AI to ask
  intelligent follow-up questions so that it simulates a real
  interviewer, not a script.
\item
  \passthrough{\lstinline!US-013!} As a user, I want to choose an
  adversarial interviewer persona so that I can practice high-pressure
  scenarios.
\item
  \passthrough{\lstinline!US-014!} As a user, I want to interrupt the AI
  when it's talking so that I can correct myself or add information.
\item
  \passthrough{\lstinline!US-015!} As a user, I want to see a live
  filler word counter so that I can self-correct in real time.
\item
  \passthrough{\lstinline!US-016!} As a user, I want the AI to probe
  when I give a vague answer so that I'm forced to be specific.
\end{itemize}

\hypertarget{scoring-reports}{%
\subsubsection{Scoring \& Reports}\label{scoring-reports}}

\begin{itemize}
\tightlist
\item
  \passthrough{\lstinline!US-020!} As a user, I want a detailed
  post-session report so that I know exactly what to improve.
\item
  \passthrough{\lstinline!US-021!} As a user, I want to replay my audio
  answers so that I can hear what the interviewer heard.
\item
  \passthrough{\lstinline!US-022!} As a user, I want AI-written example
  better answers so that I know what ``good'' looks like.
\item
  \passthrough{\lstinline!US-023!} As a user, I want per-dimension
  scores so that I can prioritize what to work on.
\end{itemize}

\hypertarget{progress}{%
\subsubsection{Progress}\label{progress}}

\begin{itemize}
\tightlist
\item
  \passthrough{\lstinline!US-030!} As a user, I want a score trend chart
  so that I can see if I'm actually improving.
\item
  \passthrough{\lstinline!US-031!} As a user, I want to see which topics
  I haven't practiced so that I don't have blind spots.
\item
  \passthrough{\lstinline!US-032!} As a user, I want weekly email
  summaries so that I stay engaged even between sessions.
\end{itemize}

\hypertarget{account}{%
\subsubsection{Account}\label{account}}

\begin{itemize}
\tightlist
\item
  \passthrough{\lstinline!US-040!} As a user, I want to sign in with
  Google so that I don't need another password.
\item
  \passthrough{\lstinline!US-041!} As a user, I want to delete my
  account and all data so that my privacy is protected.
\end{itemize}

\begin{center}\rule{0.5\linewidth}{0.5pt}\end{center}

\hypertarget{non-goals}{%
\subsection{9. Non-Goals}\label{non-goals}}

\textbf{These things will NOT be built in Phase 1. Document them
explicitly so you don't scope-creep:}

\begin{itemize}
\tightlist
\item
  ❌ Mobile app (iOS or Android)
\item
  ❌ Live interview assistance (listening to a real interview and
  whispering answers)
\item
  ❌ Video recording or video analysis
\item
  ❌ Custom voice cloning
\item
  ❌ Multi-language support (English only)
\item
  ❌ Peer-to-peer matching
\item
  ❌ Company-specific preparation packs beyond the question bank
\item
  ❌ ATS resume scoring
\item
  ❌ Coding environment (no code editor/execution)
\item
  ❌ Integration with job boards or ATS systems
\item
  ❌ Team/enterprise accounts
\item
  ❌ Offline/on-device processing
\end{itemize}

\begin{center}\rule{0.5\linewidth}{0.5pt}\end{center}

\hypertarget{success-metrics-kpis}{%
\subsection{10. Success Metrics \& KPIs}\label{success-metrics-kpis}}

\hypertarget{phase-1-launch-criteria-before-calling-it-launched}{%
\subsubsection{Phase 1 Launch Criteria (before calling it
``launched'')}\label{phase-1-launch-criteria-before-calling-it-launched}}

\begin{itemize}
\tightlist
\item[$\square$]
  End-to-end voice latency \textless{} 2 seconds for 95\% of turns
\item[$\square$]
  ASR word error rate \textless{} 10\% on clean English speech
\item[$\square$]
  Session completion rate \textgreater{} 50\% (user reaches end of
  planned session)
\item[$\square$]
  Post-session report generated in \textless{} 30 seconds
\item[$\square$]
  Zero data loss on session transcript
\item[$\square$]
  Uptime \textgreater{} 99\% over 7-day rolling window
\end{itemize}

\hypertarget{product-health-metrics-post-launch-weekly-tracking}{%
\subsubsection{Product Health Metrics (post-launch, weekly
tracking)}\label{product-health-metrics-post-launch-weekly-tracking}}

\begin{longtable}[]{@{}llll@{}}
\toprule\noalign{}
Metric & Week 4 Target & Week 12 Target & Why It Matters \\
\midrule\noalign{}
\endhead
\bottomrule\noalign{}
\endlastfoot
DAU & 10 & 50 & Core engagement signal \\
DAU/MAU & 15\% & 20\% & Retention quality \\
Sessions per user per week & 2 & 3 & Habit formation \\
Avg session length & 15 min & 20 min & Depth of engagement \\
D1 Retention & 20\% & 25\% & First impression quality \\
D7 Retention & 8\% & 12\% & Product stickiness \\
Free-to-Pro conversion & 1\% & 4\% & Monetization viability \\
NPS & - & \textgreater{} 30 & Word of mouth potential \\
Session completion rate & 55\% & 70\% & Product quality \\
\end{longtable}

\hypertarget{business-metrics-post-monetization}{%
\subsubsection{Business Metrics
(post-monetization)}\label{business-metrics-post-monetization}}

\begin{itemize}
\tightlist
\item
  Monthly Recurring Revenue (MRR)
\item
  Customer Acquisition Cost (CAC) --- target \textless{} 3x monthly
  revenue
\item
  Lifetime Value (LTV) --- target \textgreater{} 6 months retention
\item
  Churn rate --- target \textless{} 8\%/month
\end{itemize}

\begin{center}\rule{0.5\linewidth}{0.5pt}\end{center}

\hypertarget{assumptions-vs.-decisions}{%
\subsection{11. Assumptions
vs.~Decisions}\label{assumptions-vs.-decisions}}

This section is critical. It separates what we've \emph{decided} from
what we're \emph{assuming to be true}. Assumptions must be validated and
can change. Decisions are architectural commitments.

\hypertarget{assumptions-may-change-validate-as-you-build}{%
\subsubsection{ASSUMPTIONS (may change --- validate as you
build)}\label{assumptions-may-change-validate-as-you-build}}

\begin{longtable}[]{@{}
  >{\raggedright\arraybackslash}p{(\columnwidth - 4\tabcolsep) * \real{0.3333}}
  >{\raggedright\arraybackslash}p{(\columnwidth - 4\tabcolsep) * \real{0.3333}}
  >{\raggedright\arraybackslash}p{(\columnwidth - 4\tabcolsep) * \real{0.3333}}@{}}
\toprule\noalign{}
\begin{minipage}[b]{\linewidth}\raggedright
Assumption
\end{minipage} & \begin{minipage}[b]{\linewidth}\raggedright
Risk if Wrong
\end{minipage} & \begin{minipage}[b]{\linewidth}\raggedright
How to Validate
\end{minipage} \\
\midrule\noalign{}
\endhead
\bottomrule\noalign{}
\endlastfoot
Groq free tier remains usable at our scale & Groq throttles us → need
paid plan & Monitor rate limit hits weekly \\
Oracle Cloud ARM capacity available in our region & Can't provision →
delayed launch & Provision in 2 regions simultaneously \\
Kokoro TTS quality is sufficient for user retention & Users find voice
robotic → use Cartesia & User interview, NPS survey after 50 sessions \\
Browser audio APIs are sufficient for voice capture & Mic quality too
low for ASR & Test on Chrome/Firefox/Safari/Edge on Day 1 \\
Users prefer push-to-talk over always-on & Users find push-to-talk
cumbersome & A/B test both in Phase 1 \\
English-only is acceptable for launch user base & International users
churn → limit TAM & Survey first 100 users on language needs \\
2-second end-to-end latency is acceptable & Users find it too slow →
kills conversational feel & User test with explicit latency
measurement \\
Resume parsing via LLM is accurate enough & Hallucinated resume data
corrupts questions & Manual audit of first 50 parsed resumes \\
\end{longtable}

\hypertarget{decisions-architectural-commitments-change-requires-a-new-adr}{%
\subsubsection{DECISIONS (architectural commitments --- change requires
a new
ADR)}\label{decisions-architectural-commitments-change-requires-a-new-adr}}

\begin{longtable}[]{@{}
  >{\raggedright\arraybackslash}p{(\columnwidth - 4\tabcolsep) * \real{0.3333}}
  >{\raggedright\arraybackslash}p{(\columnwidth - 4\tabcolsep) * \real{0.3333}}
  >{\raggedright\arraybackslash}p{(\columnwidth - 4\tabcolsep) * \real{0.3333}}@{}}
\toprule\noalign{}
\begin{minipage}[b]{\linewidth}\raggedright
Decision
\end{minipage} & \begin{minipage}[b]{\linewidth}\raggedright
Rationale
\end{minipage} & \begin{minipage}[b]{\linewidth}\raggedright
ADR Reference
\end{minipage} \\
\midrule\noalign{}
\endhead
\bottomrule\noalign{}
\endlastfoot
WebSocket-based real-time architecture & Only option for bidirectional
audio streaming & ADR-001 \\
Web-first (no mobile in Phase 1) & Mobile audio APIs add 4--6 weeks of
bugs & ADR-002 \\
Hosted LLM providers (not self-hosted) & Self-hosted LLMs too slow on
CPU for real-time voice & ADR-003 \\
Groq as primary LLM & Sub-100ms TTFT, free tier, OpenAI-compatible API &
ADR-004 \\
Kokoro TTS (self-hosted) & Apache 2.0, CPU-capable, unlimited usage,
8.5/10 quality & ADR-005 \\
Deepgram for streaming ASR & Best latency, \$200 free credit, real-time
WebSocket native & ADR-006 \\
Supabase for auth + PostgreSQL & Bundles auth, saves weeks of work in
Phase 1 & ADR-007 \\
Valkey for session cache & Open-source Redis fork, self-hosted on Oracle
ARM & ADR-008 \\
Oracle Cloud ARM as compute & 24 GB RAM, always free, enough for entire
stack & ADR-009 \\
Structured latency logging from Day 1 & Can't optimize what you can't
measure & ADR-010 \\
\end{longtable}

\begin{center}\rule{0.5\linewidth}{0.5pt}\end{center}

\hypertarget{risks}{%
\subsection{12. Risks}\label{risks}}

\hypertarget{technical-risks}{%
\subsubsection{Technical Risks}\label{technical-risks}}

\textbf{R-001: Real-time latency target is hard to achieve.} Getting
end-to-end \textless{} 2 seconds is genuinely difficult. It requires
every pipeline stage to perform near-optimally. \emph{Mitigation:}
Instrument from Day 1. Have a latency budget per stage. Accept 3 seconds
for Phase 1 launch, optimize to 2 seconds in Phase 1.5.

\textbf{R-002: ASR accuracy degrades in real conditions.} WER of 5\% in
benchmarks becomes 10--15\% with background noise, accents, and
technical vocabulary. \emph{Mitigation:} Deepgram Nova-3 is the most
robust streaming model. Add Whisper post-processing for technical terms.

\textbf{R-003: Oracle Cloud ARM availability.} Oracle may reclaim
instances idle \textless{} 20\% CPU or ``Out of Capacity'' on ARM
provisioning. \emph{Mitigation:} Upgrade to Pay-As-You-Go (free, just
needs credit card) for better availability. Use keep-alive script to
maintain CPU activity.

\textbf{R-004: Groq rate limits hit at scale.} Free tier is 30 RPM and
500K tokens/day per model. At 100 concurrent users, this runs out in
minutes. \emph{Mitigation:} Multi-provider fallback chain. Monitor rate
limit hits. Budget \$20/month for paid Groq when needed --- much cheaper
than OpenAI.

\hypertarget{product-risks}{%
\subsubsection{Product Risks}\label{product-risks}}

\textbf{R-005: Product feels robotic, not like a real interviewer.} If
the AI responses are generic, users won't return. \emph{Mitigation:}
Invest heavily in system prompt design. A/B test prompt variations. Get
10 real users to give honest feedback before calling it launched.

\textbf{R-006: Low retention.} Career tools have natural ``positive
churn'' (user gets job, cancels). \emph{Mitigation:} Build for referrals
(``I used SpeakPrep and got the job --- my friend is looking\ldots{}'').
Track cohort-based retention, not just raw numbers.

\textbf{R-007: High free-to-paid friction.} EdTech freemium conversion
averages only 2.6\%. \emph{Mitigation:} Make free tier genuinely useful
but leave meaningful value in Pro (full reports, audio playback,
unlimited sessions). Don't cripple free tier to the point of being
useless.

\begin{center}\rule{0.5\linewidth}{0.5pt}\end{center}

\hypertarget{monetization-strategy}{%
\subsection{13. Monetization Strategy}\label{monetization-strategy}}

\hypertarget{pricing-phase-1}{%
\subsubsection{Pricing (Phase 1)}\label{pricing-phase-1}}

\textbf{Free Tier --- ``Try it''} - 3 sessions/week, 15 minutes each -
Basic post-session score (overall + top-level dimensions) - No audio
playback, no full transcript - 7-day history

\textbf{Pro --- \$12/month or \$89/year (save 38\%)} - Unlimited
sessions, up to 60 minutes each - Full post-session reports (all 5
dimensions, per-question breakdown) - Audio playback + full transcript -
90-day progress history + trend charts - All interviewer personas
(including adversarial) - Company-specific question sets -
Resume-personalized questions - Priority support

\hypertarget{why-12month}{%
\subsubsection{Why \$12/month}\label{why-12month}}

\begin{itemize}
\tightlist
\item
  Below competitor range (\$39--\$148/month)
\item
  Cheaper than one hour of human coaching (\$150+)
\item
  Comparable to a gym membership --- users understand recurring
  fitness/skills investment
\item
  Low enough to not require approval from a partner/parent
\item
  High enough to be a real business at scale: 1,000 Pro users = \$12,000
  MRR
\end{itemize}

\hypertarget{revenue-model-sustainability}{%
\subsubsection{Revenue Model
Sustainability}\label{revenue-model-sustainability}}

At 1\% free-to-paid conversion with 1,000 active free users: 10 paid
users × \$12 = \$120 MRR. At 4\% conversion with 5,000 free users: 200
paid × \$12 = \$2,400 MRR. At 4\% conversion with 25,000 free users:
1,000 paid × \$12 = \$12,000 MRR (product-market fit territory).

\textbf{Variable cost per Pro user per month:} - Deepgram ASR:
\textasciitilde\$2 (unlimited sessions, avg 4 sessions/week × 30 min ×
\$0.0077/min) - Groq LLM: \textasciitilde\$0.50 (avg 200 turns × 500
tokens × \$0.0008/1K tokens) - Infrastructure: \textasciitilde\$0.01
(Oracle ARM is free; amortize domain/monitoring) - \textbf{Total
variable cost: \textasciitilde\$2.50/user/month} - \textbf{Gross margin:
\textasciitilde79\%} ✓

\begin{center}\rule{0.5\linewidth}{0.5pt}\end{center}

\hypertarget{competitive-differentiation}{%
\subsection{14. Competitive
Differentiation}\label{competitive-differentiation}}

\hypertarget{why-speakprep-wins-or-at-least-earns-its-users}{%
\subsubsection{Why SpeakPrep Wins (or at least earns its
users)}\label{why-speakprep-wins-or-at-least-earns-its-users}}

\begin{longtable}[]{@{}
  >{\raggedright\arraybackslash}p{(\columnwidth - 10\tabcolsep) * \real{0.1667}}
  >{\raggedright\arraybackslash}p{(\columnwidth - 10\tabcolsep) * \real{0.1667}}
  >{\raggedright\arraybackslash}p{(\columnwidth - 10\tabcolsep) * \real{0.1667}}
  >{\raggedright\arraybackslash}p{(\columnwidth - 10\tabcolsep) * \real{0.1667}}
  >{\raggedright\arraybackslash}p{(\columnwidth - 10\tabcolsep) * \real{0.1667}}
  >{\raggedright\arraybackslash}p{(\columnwidth - 10\tabcolsep) * \real{0.1667}}@{}}
\toprule\noalign{}
\begin{minipage}[b]{\linewidth}\raggedright
Dimension
\end{minipage} & \begin{minipage}[b]{\linewidth}\raggedright
SpeakPrep
\end{minipage} & \begin{minipage}[b]{\linewidth}\raggedright
ChatGPT/Claude
\end{minipage} & \begin{minipage}[b]{\linewidth}\raggedright
Yoodli
\end{minipage} & \begin{minipage}[b]{\linewidth}\raggedright
Final Round AI
\end{minipage} & \begin{minipage}[b]{\linewidth}\raggedright
Big Interview
\end{minipage} \\
\midrule\noalign{}
\endhead
\bottomrule\noalign{}
\endlastfoot
Real-time voice conversation & ✅ & ❌ (text) & ✅ & ✅ & ❌ \\
Intelligent probing/follow-ups & ✅ & ❌ & ❌ & Partial & ❌ \\
Resume-personalized questions & ✅ & Manual setup & ❌ & ❌ & ❌ \\
Adaptive difficulty & ✅ & ❌ & ❌ & ❌ & ❌ \\
Progress tracking across sessions & ✅ & ❌ & Basic & ❌ & Basic \\
Detailed AI scoring rubric & ✅ & Variable & Speech only & ❌ &
Generic \\
End-to-end interview coverage & ✅ & Ad hoc & Delivery only & Behavioral
& Limited \\
Price & \$12/mo & \$20/mo (Plus) & \$10--20/mo & \$96--148/mo &
\$39--299 \\
\end{longtable}

\hypertarget{the-one-sentence-differentiator}{%
\subsubsection{The One-Sentence
Differentiator}\label{the-one-sentence-differentiator}}

\begin{quote}
``SpeakPrep is the only voice AI coach that asks intelligent follow-up
questions, adapts difficulty to your skill level, and shows you exactly
how you're improving over time --- for less than the cost of one hour
with a human coach.''
\end{quote}

\begin{center}\rule{0.5\linewidth}{0.5pt}\end{center}

\emph{End of Document 1 --- PRD + Product} \emph{Next: Document 2 ---
System Design + Architecture} \emph{Cross-reference: All feature IDs
(F-001 through F-010) and user stories (US-001 through US-041) are
referenced in Doc 2 and Doc 3.}

\end{document}
